\section{Implementation de l'espace administrateur}

Cette section presente les principales interfaces reservees a l'administrateur.
Chaque page repond a un besoin de pilotage precis : suivre l'activite du centre,
traiter les inscriptions, gerer les etudiants, organiser les formations et
controler les paiements. Les captures d'ecran illustrent les fonctions visibles
par l'utilisateur final, tandis que les paragraphes suivants resument leur role
dans le fonctionnement global de la plateforme.

\subsection{Tableau de bord}

Le tableau de bord administrateur centralise les indicateurs essentiels du centre
Info Saloua. Il permet d'avoir une vision rapide sur le nombre d'etudiants, les
inscriptions actives, les nouvelles demandes, les revenus et le taux de
conversion. Cette page sert donc de point d'entree pour surveiller l'activite
globale et identifier rapidement les elements qui demandent une action.

\begin{table}[H]
	\centering
	\renewcommand{\arraystretch}{1.5}
	\begin{tabular}{>{\columncolor{tableheader}\bfseries\color{white}}m{5cm}
			>{\columncolor{tablerowA}}m{8.5cm}}
		\rowcolor{tableheader}
		\textcolor{white}{\textbf{Indicateur}} &
		\textcolor{white}{\textbf{Description}} \\
		Etudiants totaux       & Nombre total d'etudiants actifs \\
		Inscriptions actives   & Nombre d'inscriptions en cours \\
		Nouveaux ce mois       & Nouveaux etudiants du mois \\
		Revenu mensuel         & Chiffre d'affaires du mois en cours \\
		Montant total encaisse & Cumul des paiements recus \\
		Demandes en attente    & Demandes non encore traitees \\
		Taux de conversion     & Ratio demandes approuvees / demandes totales \\
	\end{tabular}
	\caption{Indicateurs du tableau de bord administrateur}
\end{table}

\begin{figure}[H]
	\centering
	\includegraphics[width=0.95\textwidth]{images/admin_dashboard.png}
	\caption{Tableau de bord de l'espace administrateur}
	\label{fig:admin_dashboard}
\end{figure}

\subsection{Gestion des demandes d'inscription}

La page des demandes d'inscription regroupe toutes les candidatures envoyees
depuis les formulaires publics. L'administrateur peut consulter les informations
du candidat, verifier la formation demandee, le montant associe, le mode de
paiement et le statut du dossier. Les actions d'approbation ou de rejet
permettent de transformer une demande validee en inscription effective, tout en
gardant une trace claire du traitement realise.

\begin{figure}[H]
	\centering
	\includegraphics[width=0.95\textwidth]{images/admin_demandes.png}
	\caption{Interface de gestion des demandes d'inscription}
	\label{fig:admin_demandes}
\end{figure}

\subsection{Gestion des etudiants}

Cette interface assure la gestion des comptes etudiants. Elle permet de creer,
modifier ou supprimer un profil, de consulter les coordonnees, de suivre le
statut du compte et de mettre a jour les informations personnelles. Elle est
importante pour maintenir une base d'etudiants fiable et exploitable par les
autres modules, notamment les inscriptions, paiements, notifications et
plannings.

\begin{figure}[H]
	\centering
	\includegraphics[width=0.95\textwidth]{images/admin_etudiants.png}
	\caption{Interface de gestion des etudiants}
	\label{fig:admin_etudiants}
\end{figure}

\subsection{Gestion des formations}

La page des formations permet d'administrer le catalogue propose par le centre.
Chaque formation contient des informations telles que le titre, la categorie, le
niveau, la description, le prix, la capacite et le statut. L'administrateur peut
ainsi garder le catalogue a jour et preparer les formations qui seront visibles
dans les parcours d'inscription.

\begin{figure}[H]
	\centering
	\includegraphics[width=0.95\textwidth]{images/admin_formations.png}
	\caption{Interface de gestion des formations}
	\label{fig:admin_formations}
\end{figure}

\subsection{Gestion du planning}

Le module de planning organise les seances liees aux formations. Il permet de
definir la date, l'heure, la salle, le groupe et le statut d'une seance. Cette
page facilite la coordination entre l'administration et les etudiants, car les
seances programmees sont ensuite reutilisees dans l'espace etudiant pour afficher
le calendrier personnel et la prochaine seance.

\begin{figure}[H]
	\centering
	\includegraphics[width=0.95\textwidth]{images/admin_planning.png}
	\caption{Interface de gestion du planning des seances}
	\label{fig:admin_planning}
\end{figure}

\subsection{Gestion des paiements}

La page des paiements donne une vue d'ensemble sur les transactions associees aux
inscriptions. Elle affiche les references, les etudiants, les montants, les
methodes de paiement et les statuts. Cette interface aide l'administrateur a
verifier les paiements regles, reperer les paiements en attente et suivre le
revenu genere par les formations.

\begin{figure}[H]
	\centering
	\includegraphics[width=0.95\textwidth]{images/admin_paiements.png}
	\caption{Interface de gestion des paiements}
	\label{fig:admin_paiements}
\end{figure}

\subsection{Notifications administrateur}

La page des notifications regroupe les alertes utiles pour l'administration :
nouvelles demandes, paiements, evenements de planning ou informations systeme.
Elle permet de consulter les messages importants et de marquer les notifications
comme lues. Ce module renforce le suivi operationnel en evitant que les actions
urgentes restent dispersees entre plusieurs pages.

\begin{figure}[H]
	\centering
	\includegraphics[width=0.95\textwidth]{images/admin_notifications.png}
	\caption{Centre de notifications de l'espace administrateur}
	\label{fig:admin_notifications}
\end{figure}

\subsection{Rapports statistiques}

La page des rapports consolide les donnees academiques et financieres sur une
annee selectionnee. Elle presente le revenu annuel, la croissance, les
inscriptions, le taux de validation, la repartition des formations, l'etat des
paiements et l'historique des performances. Ces visualisations aident
l'administration a prendre des decisions basees sur des donnees reelles.

\begin{figure}[H]
	\centering
	\includegraphics[width=0.95\textwidth]{images/admin_rapports.png}
	\caption{Page des rapports statistiques de l'espace administrateur}
	\label{fig:admin_rapports}
\end{figure}

\section{Implementation de l'espace etudiant}

L'espace etudiant regroupe les informations personnelles de l'apprenant apres
connexion. Il permet de consulter les formations suivies, le planning, les
paiements et les notifications. L'objectif est de donner a chaque etudiant une
vue claire de son parcours sans passer par l'administration.

\subsection{Tableau de bord etudiant}

Le tableau de bord etudiant resume les informations les plus importantes :
formations actives, prochaines seances, paiements recents et notifications. Il
affiche aussi une salutation personnalisee et la date du jour. Cette page sert de
point d'accueil apres la connexion et oriente l'etudiant vers les actions
principales.

\begin{figure}[H]
	\centering
	\includegraphics[width=0.95\textwidth]{images/etudiant_dashboard.png}
	\caption{Tableau de bord de l'espace etudiant}
	\label{fig:etudiant_dashboard}
\end{figure}

\subsection{Mes formations}

La page Mes formations presente les formations auxquelles l'etudiant est inscrit.
Chaque carte affiche le titre, la categorie, la description et l'etat
d'avancement. Cette interface permet a l'etudiant de suivre son parcours et de
retrouver rapidement les informations liees a ses inscriptions.

\begin{figure}[H]
	\centering
	\includegraphics[width=0.92\textwidth]{images/etudiant_formations.png}
	\caption{Page Mes formations de l'espace etudiant}
	\label{fig:etudiant_formations}
\end{figure}

\subsection{Planning}

Le planning etudiant affiche les seances programmees pour les formations suivies.
Il precise les horaires, les salles, les groupes et les informations utiles pour
assister aux cours. Cette page est directement liee au planning admin, ce qui
garantit que les changements effectues par l'administration sont repercutes dans
l'espace etudiant.

\begin{figure}[H]
	\centering
	\includegraphics[width=0.92\textwidth]{images/etudiant_planning.png}
	\caption{Page planning de l'espace etudiant}
	\label{fig:etudiant_planning}
\end{figure}

\subsection{Paiements}

La page Paiements permet a l'etudiant de consulter l'historique de ses
transactions. Elle affiche les references, les montants, les methodes de paiement
et les statuts. L'etudiant peut ainsi verifier les paiements deja regles et
identifier ceux qui restent en attente.

\begin{figure}[H]
	\centering
	\includegraphics[width=0.92\textwidth]{images/etudiant_paiements.png}
	\caption{Page des paiements de l'espace etudiant}
	\label{fig:etudiant_paiements}
\end{figure}

\subsection{Notifications etudiant}

Les notifications etudiant regroupent les messages lies au parcours de
l'apprenant : validation d'inscription, rappel de paiement, information de
planning ou autre annonce importante. L'etudiant peut consulter les notifications
recentes et les marquer comme lues, ce qui facilite le suivi des informations
envoyees par l'administration.

\begin{figure}[H]
	\centering
	\includegraphics[width=0.92\textwidth]{images/etudiant_notifications.png}
	\caption{Centre de notifications de l'espace etudiant}
	\label{fig:etudiant_notifications}
\end{figure}

\subsection{Profil utilisateur}

La page Profil est commune aux utilisateurs connectes. Elle permet de modifier
les informations personnelles, de mettre a jour le mot de passe, de gerer
l'avatar et, si necessaire, de supprimer le compte. Cette page complete les
espaces admin et etudiant en donnant a chaque utilisateur le controle sur ses
donnees de connexion et son identite visuelle.

\begin{figure}[H]
	\centering
	\includegraphics[width=0.92\textwidth]{images/profil_utilisateur.png}
	\caption{Page de modification du profil utilisateur}
	\label{fig:profil_utilisateur}
\end{figure}

\section{Deploiement sur Railway}

Railway est utilise pour heberger l'application et la base de donnees PostgreSQL
en environnement de production. Le deploiement repose sur les variables
d'environnement configurees depuis le tableau de bord Railway, ce qui permet de
separer les informations sensibles du code source. La commande de deploiement
execute les migrations, puis regenere les caches Laravel afin de preparer
l'application pour la production.

\begin{figure}[H]
	\centering
	\includegraphics[width=0.92\textwidth]{images/railway_dashboard.png}
	\caption{Dashboard Railway : deploiement de l'application et de PostgreSQL}
	\label{fig:railway}
\end{figure}

\section{Conclusion}

Ce chapitre a presente la realisation technique de la plateforme Info Saloua,
depuis les espaces prives jusqu'au deploiement. Les interfaces administrateur
assurent le pilotage du centre, tandis que l'espace etudiant facilite le suivi
des formations, seances, paiements et notifications. L'ensemble forme une
solution coherente, connectee a la base de donnees et prete a etre exploitee en
environnement de production.
